% =================================================================================
% Response-to-reviewers letter template (Elsevier journals).
% Letter-style layout (article class), compiled with pdfLaTeX -- Editorial
% Manager's default engine:
%
%     latexmk -pdf main.tex
%
% Placeholder convention: every value that belongs to one specific paper is
% written as {{FIELD-NAME}} -- double curly braces, ALL-CAPS, hyphen-separated.
% The braces are LaTeX grouping characters, so an unsubstituted template still
% compiles and the placeholder shows up in the PDF where it will be replaced.
% Find them all with:
%
%     grep -ron '{{[A-Z0-9-]\+}}' .
%
% See README.md for the full placeholder list.
% =================================================================================

\documentclass[a4paper,twoside,11pt]{article}

% --- Encoding, Language and Fonts (pdfLaTeX-compatible) ------------------------
% Editorial Manager defaults to pdfLaTeX.  We avoid fontspec / system fonts and
% use the standard Helvetica-compatible sans-serif family, which preserves the
% Arial/Helvetica visual style without requiring local Arial.
\usepackage[english]{babel}
\usepackage[utf8]{inputenc}
\usepackage[T1]{fontenc}
\usepackage{textcomp}
\usepackage[scaled=0.95]{helvet}
\renewcommand{\familydefault}{\sfdefault}

% --- Page Geometry --------------------------------------------------------------
\usepackage[
    a4paper,
    top=2.5cm,
    bottom=2.5cm,
    left=1.8cm,
    right=1.8cm,
    headheight=14pt
]{geometry}

% --- Typography -----------------------------------------------------------------
\usepackage{setspace}
\setstretch{1}
\usepackage{calc}
\usepackage{verbatim}
\usepackage{url}
\usepackage{csquotes}

% --- Math -----------------------------------------------------------------------
\usepackage{amsmath}
\usepackage{amssymb}

% --- Graphics & Colors ----------------------------------------------------------
\usepackage[dvipsnames]{xcolor}

% --- Lists ----------------------------------------------------------------------
\usepackage{enumitem}

% --- Header & Footer ------------------------------------------------------------
\usepackage{fancyhdr}
\pagestyle{fancy}
\fancyhf{}
\fancyhead[LO]{\leftmark}
\fancyhead[RE]{\myJournal{} -- \myShortTitle}
\fancyhead[RO,LE]{\thepage}
\fancyfoot[C]{\thepage}
\renewcommand{\headrulewidth}{0pt}
\renewcommand{\footrulewidth}{0pt}

% --- Hyperlinks (loaded before tcolorbox so the two cooperate) ------------------
\usepackage[
    bookmarks, bookmarksopen=true, bookmarksnumbered=true,
    colorlinks=true, linkcolor=black,
    citecolor=blue!50!black, urlcolor=blue!50!black,
    linktoc=all
]{hyperref}

% --- Author Affiliations --------------------------------------------------------
\usepackage{authblk}
\renewcommand\Affilfont{\itshape\small}

% --- Bibliography ---------------------------------------------------------------
% No bibliography package is loaded on purpose: literature references inside a
% response letter are given inline in author-year plain text (e.g.\ ``Guo
% et~al., CVPR 2019, ...''), so the reader can cross-check them against the
% manuscript's own bibliography without an extra bib pass.  If you do need a
% real bibliography, add natbib/biblatex here -- nothing else in the template
% depends on its absence.

% --- tcolorbox (visual engine) --------------------------------------------------
\usepackage[most]{tcolorbox}
\tcbuselibrary{skins, breakable}

% --- Pantone-inspired editorial palette -----------------------------------------
% Coated (C) hex equivalents drawn from the Pantone Solid Coated guide; tuned
% for legible PDF rendering.  Three families (gray / blue / green) replace the
% default xcolor names {gray!60!black, cyan!70!black, green!45!black} so that
% the letter reads as a single coherent palette.
%
% Pantone Cool Gray 10 C   -> reviewer / editor comment frame
% Pantone Cool Gray 1 C    -> comment title fill
\definecolor{pantonegray}{HTML}{4F5152}
\definecolor{pantonegraytitle}{HTML}{ECEDED}
% Pantone 295 C            -> our-response frame (deep editorial blue)
% Pantone 657 C            -> our-response title fill
\definecolor{pantoneblue}{HTML}{1F3A68}
\definecolor{pantonebluetitle}{HTML}{D9E2F1}
% Pantone 7484 C           -> revised-manuscript-text frame (forest green)
% Pantone 558 C            -> revised-text title fill
% Pantone 559 C 5\% tint   -> revised-text body wash
\definecolor{pantonegreen}{HTML}{2D6A4F}
\definecolor{pantonegreentitle}{HTML}{D8E8DC}
\definecolor{pantonegreenbody}{HTML}{F5F9F6}

% TOC entries use ``linktoc=all'' from the hyperref load, so re-pointing
% linkcolor at Pantone blue paints the entire table of contents (entry
% text + page numbers + dot leaders' clickable region) in the same blue
% used for citations and URLs.  This is what gives the letter its
% editorial-blue TOC without pulling in the elsarticle layout.
\hypersetup{
    linkcolor=pantoneblue,
    citecolor=pantoneblue,
    urlcolor=pantoneblue
}

% --- Shared base style: title bar INSIDE the box -------------------------------
% An earlier revision of this template used ``attach boxed title to top
% left=...'', which let each box's title float above the frame; that worked
% aesthetically but the floating title was not reserved in the surrounding
% paragraph flow, so on dense pages the next box's title clipped into the
% previous box's bottom frame.  We now place the title bar INSIDE the frame.
% Because the title is part of the box's vertical extent, it cannot collide
% with the previous box no matter how a page break or paragraph descender
% shifts the layout.
\tcbset{
    letterboxbase/.style={
        enhanced, breakable,
        boxrule=0.5pt,
        arc=2pt,
        left=10pt, right=10pt, top=4pt, bottom=8pt,
        boxsep=2pt,
        before skip=12pt, after skip=8pt,
        fonttitle=\sffamily\bfseries\small,
        coltitle=black,
        toptitle=3pt, bottomtitle=3pt,
        titlerule=0pt
    }
}

% Loosen line-breaking for dense bullet items with inline math/code.
\setlength{\emergencystretch}{2em}
\tolerance=2500


% =================================================================================
% Metadata -- the six per-paper commands.  Substitute the {{...}} placeholders
% and nothing else in this block needs to change.
% =================================================================================
\newcommand{\myShortTitle}{Response to reviewers}
\newcommand{\myJournal}{{{JOURNAL-NAME}}}
\newcommand{\myManuscriptID}{{{MANUSCRIPT-ID}}}
\newcommand{\myManuscriptTitle}{{{MANUSCRIPT-TITLE}}}
\newcommand{\myEditorName}{{{EDITOR-NAME}}}
\newcommand{\myRevisionRound}{{{REVISION-ROUND}}}

\newcommand{\horizontalline}{\noindent\rule{\textwidth}{0.6pt}}


% =================================================================================
% Counters and environments for editor / reviewer / revised-text boxes.
%
% NUMBERING IS BY ORDER OF APPEARANCE.  Both comment environments call
% \stepcounter, so ``Comment 3'' is simply the third reviewercomment in the
% file -- there is no explicit number to pass.  Reordering the environments
% renumbers the letter.  \startReviewerSection resets the reviewer counter,
% so every reviewer starts again at 1.
% =================================================================================

% --- Editor counters / names ----------------------------------------------------
\newcounter{editorcomment}
\newcommand{\currenteditorcommentfulltitle}{}
\newcommand{\currenteditorcommentshorttitle}{}

% --- Reviewer counters / names --------------------------------------------------
\newcounter{reviewercomment}
\newcommand{\currentreviewername}{}
\newcommand{\currentreviewercommentfulltitle}{}
\newcommand{\currentreviewercommentshorttitle}{}

\newcommand{\startReviewerSection}[1]{%
    \clearpage
    \section{Response to #1}%
    \par\nobreak\vspace{1.0\baselineskip}%
    \setcounter{reviewercomment}{0}%
    \gdef\currentreviewername{#1}%
}

% --- Editor comment (Pantone Cool Gray 10) -------------------------------------
\newenvironment{editorcomment}{%
    \stepcounter{editorcomment}%
    \xdef\currenteditorcommentfulltitle{Editor Comment \theeditorcomment}%
    \xdef\currenteditorcommentshorttitle{Comment \theeditorcomment}%
    \addcontentsline{toc}{subsection}{\currenteditorcommentfulltitle}%
    \begin{tcolorbox}[
        letterboxbase,
        colback=white, colframe=pantonegray,
        colbacktitle=pantonegraytitle,
        title=\currenteditorcommentfulltitle
    ]%
    \itshape
}{%
    \end{tcolorbox}%
}

% --- Editor response (Pantone 295 C) -------------------------------------------
\newenvironment{editorresponse}{%
    \addcontentsline{toc}{subsubsection}{Response to \currenteditorcommentshorttitle}%
    \begin{tcolorbox}[
        letterboxbase,
        colback=white, colframe=pantoneblue,
        colbacktitle=pantonebluetitle,
        title=Response to \currenteditorcommentshorttitle
    ]%
}{%
    \end{tcolorbox}%
}

% --- Reviewer comment (Pantone Cool Gray 10) -----------------------------------
\newenvironment{reviewercomment}{%
    \stepcounter{reviewercomment}%
    \xdef\currentreviewercommentfulltitle{\currentreviewername\ Comment \thereviewercomment}%
    \xdef\currentreviewercommentshorttitle{Comment \thereviewercomment}%
    \addcontentsline{toc}{subsection}{\currentreviewercommentfulltitle}%
    \begin{tcolorbox}[
        letterboxbase,
        colback=white, colframe=pantonegray,
        colbacktitle=pantonegraytitle,
        title=\currentreviewercommentfulltitle
    ]%
    \itshape
}{%
    \end{tcolorbox}%
}

% --- Reviewer response (Pantone 295 C) -----------------------------------------
\newenvironment{reviewerresponse}{%
    \addcontentsline{toc}{subsubsection}{Response to \currentreviewercommentshorttitle}%
    \begin{tcolorbox}[
        letterboxbase,
        colback=white, colframe=pantoneblue,
        colbacktitle=pantonebluetitle,
        title=Response to \currentreviewercommentshorttitle
    ]%
}{%
    \end{tcolorbox}%
}

% --- Revised manuscript text (Pantone 7484 C) ----------------------------------
% Used inside a \begin{reviewerresponse} to display the exact wording that
% has been added to the manuscript.
\newenvironment{changes}{%
    \begin{tcolorbox}[
        letterboxbase,
        colback=pantonegreenbody, colframe=pantonegreen,
        colbacktitle=pantonegreentitle,
        title=Revised manuscript text
    ]%
}{%
    \end{tcolorbox}%
}


% =================================================================================
% Author block
% =================================================================================
% authblk: repeat \author/\affil pairs for co-authors.  The ``*'' marks the
% corresponding author; the matching \affil is the corresponding-author note.
\author[a,*]{{{AUTHOR-NAME}}}
\affil[a]{{{AUTHOR-AFFILIATION}}}
\date{}


\begin{document}

\thispagestyle{plain}
\title{\textbf{Response Letter to the Editor and Reviewers} \\
       \large{\myJournal} \\[0.2em]
       \normalsize Manuscript \myManuscriptID\ --- \myRevisionRound}
\maketitle

\vspace{-2em}

% --- Manuscript Information panel ---------------------------------------------
\begin{tcolorbox}[
    enhanced, breakable,
    colback=pantonegraytitle!40,
    colframe=pantonegray,
    boxrule=0.5pt, arc=2pt,
    left=10pt, right=10pt, top=4pt, bottom=8pt,
    boxsep=2pt,
    fonttitle=\sffamily\bfseries\small,
    coltitle=black,
    colbacktitle=pantonegraytitle,
    toptitle=3pt, bottomtitle=3pt,
    titlerule=0pt,
    title=Manuscript Information
]
    \begin{tabular}{@{} p{0.22\linewidth} p{0.74\linewidth} @{}}
        \textbf{Manuscript ID:}   & \myManuscriptID \\
        \textbf{Title:}           & \myManuscriptTitle \\
        \textbf{Submitted to:}    & \myJournal{} (Elsevier) \\
        \textbf{Handling Editor:} & \myEditorName \\
        \textbf{Revision round:}  & \myRevisionRound
    \end{tabular}
\end{tcolorbox}

\vspace{1em}

% --- Opening Letter -----------------------------------------------------------
% Keep this short and factual: thank, name the manuscript, say the revision is
% enclosed.  No promises, no summary of results -- the summary section below is
% where the substance goes.
\noindent Dear \myEditorName{} and the Reviewers,

\vspace{0.4em}

We thank you for handling our submission and for the constructive feedback from the review panel. Please find enclosed the revised version of our manuscript entitled \enquote{\myManuscriptTitle} (\myManuscriptID), originally submitted to \myJournal. Your comments have helped us improve the clarity, completeness, and experimental rigor of the paper.

\section*{Summary of Revisions Made}
\addcontentsline{toc}{section}{Summary of Revisions Made}

In this revision, we have addressed all points raised by {{REVIEWER-LIST}}. The main high-level changes are:

% Every item follows the same shape:
%   \textbf{<what changed>} (<where it landed in the manuscript>).
%   <one or two sentences of substance, with the concrete numbers>.
%   \emph{Addresses <comment ids>.}
%
% The trailing \emph{Addresses ...} sentence is not decoration: an item that
% cannot name the comment ids it serves is work nobody asked for, and it does
% not belong in this list.  Use \emph{Supports ...} for a change that helps a
% comment without being the direct answer to it, and \emph{Voluntary cleanup}
% for the small consistency fixes -- but keep those to a minimum.
% Delete the example below and write one item per substantive change.
\begin{itemize}[leftmargin=*, topsep=2pt, itemsep=2pt]
    \item \textbf{{{CHANGE-HEADLINE}}} (Sec.~{{SECTION}}, Table~{{TABLE}}). {{CHANGE-DESCRIPTION-WITH-CONCRETE-NUMBERS}} \emph{Addresses {{COMMENT-IDS}}.}
\end{itemize}

\section*{Point-by-Point Reply to the Reports}
\addcontentsline{toc}{section}{Point-by-Point Reply to the Reports}

Revised passages are quoted in change boxes below so that the corresponding manuscript updates are easy to locate. In this letter:
\begin{itemize}[leftmargin=*, topsep=2pt, itemsep=1pt]
    \item The comments from the Editor and the Reviewers are typeset in \textit{italic} and enclosed in {\color{pantonegray}gray boxes}.
    \item Our corresponding responses are presented in {\color{pantoneblue}blue boxes}.
    \item The exact textual additions that now appear in the manuscript are showcased in {\color{pantonegreen}green boxes}.
\end{itemize}

\vspace{1em}
\noindent With kind regards,

\vspace{0.5em}
\noindent {{AUTHOR-NAME}} (on behalf of all authors)

\clearpage
\tableofcontents
\clearpage

% =================================================================================
% Point-by-point responses -- one file per audience, under Reviewers/.
% Add one \input line per reviewer, in the order the decision letter uses.
% The order of these \input lines is the order of the sections in the PDF.
% =================================================================================

% Editor / meta review
% =================================================================================
% Response to the Editor (meta-review).
%
% NUMBERING IS BY ORDER OF APPEARANCE: each \begin{editorcomment} steps the
% editor counter, so the Nth editorcomment in this file becomes ``Editor
% Comment N''.  Move the environments and the numbers move with them.  Each
% editorcomment should be followed immediately by its editorresponse.
%
% The editor's summary is a separate spec from the reviewers' reports -- it is
% what the person who decides the outcome chose to emphasise.  Answer it on its
% own terms even when every point is repeated below.
% =================================================================================
\section{Response to the Editor}

% Paste the decision letter's editorial paragraph verbatim.  The environment
% already sets it in italic -- do not add quotation marks or \emph.
\begin{editorcomment}
{{EDITOR-COMMENT-VERBATIM}}
\end{editorcomment}

\begin{editorresponse}
We thank you for handling our submission and for the constructive feedback from the review panel. We have revised the manuscript following every comment from {{REVIEWER-LIST}}, and we summarize the most substantial changes below.

% One enumerate item per high-level change.  Same shape as the Summary of
% Revisions in main.tex: bold headline, concrete substance, then the comment
% ids the change answers.  Keep this list consistent with that one -- a
% mismatch between the editor summary and the point-by-point body is the first
% thing a re-reviewer notices.
\begin{enumerate}[leftmargin=*, topsep=2pt, itemsep=2pt]
    \item \textbf{{{CHANGE-HEADLINE}}} {{CHANGE-DESCRIPTION-WITH-CONCRETE-NUMBERS}} Addresses {{COMMENT-IDS}}.

    \item \textbf{{{CHANGE-HEADLINE}}} {{CHANGE-DESCRIPTION-WITH-CONCRETE-NUMBERS}} Addresses {{COMMENT-IDS}}.
\end{enumerate}

We detail each comment and our response in the following sections, preserving the original reviewer numbering from the decision letter.
\end{editorresponse}


% Reviewer 1
% =================================================================================
% Reviewer 1 -- one file per reviewer, \input from main.tex.
%
% NUMBERING IS BY ORDER OF APPEARANCE: each \begin{reviewercomment} calls
% \stepcounter, so the Nth reviewercomment in this file becomes ``Reviewer \#1
% Comment N''.  There is no number to pass and no number to keep in sync --
% but reordering the environments renumbers the letter, and the ids you cite
% elsewhere (R1-Q3, the Summary of Revisions, the HTML board) will no longer
% match.  Keep the environments in the order the decision letter uses.
%
% \startReviewerSection resets the counter, so the next reviewer file starts
% again at Comment 1.
%
% Every numbered comment gets its own reviewercomment/reviewerresponse pair,
% including the trivial ones (a typo, a missing citation).  A comment with no
% entry reads as a comment that was ignored.
% =================================================================================
\startReviewerSection{Reviewer \#1}

% --- R1-Q1: <one-line summary of the ask, for navigation> ----------------------
% Paste the reviewer's words verbatim -- do not paraphrase, do not fix their
% typos, do not merge two of their points into one box.  The environment sets
% the text in italic already.
\begin{reviewercomment}
{{REVIEWER-COMMENT-VERBATIM}}
\end{reviewercomment}

% Response order: state the action taken -> show the evidence -> point to the
% location in the revised manuscript.  Report only work that exists.
\begin{reviewerresponse}
We thank the reviewer for this observation. {{ACTION-TAKEN}} The revised text now appears in Sec.~{{SECTION}} of the manuscript, and {{EVIDENCE-OR-NUMBERS}}.
\end{reviewerresponse}

% --- R1-Q2: <one-line summary of the ask> --------------------------------------
\begin{reviewercomment}
{{REVIEWER-COMMENT-VERBATIM}}
\end{reviewercomment}

\begin{reviewerresponse}
We agree that the original description was too compact. {{ACTION-TAKEN}} The new text reads:
% The changes box quotes the exact wording that is now in the manuscript --
% copy it from the revised source, do not re-write it here.  Anything inside
% this box must be findable verbatim in the attached manuscript.
\begin{changes}
{{VERBATIM-REVISED-MANUSCRIPT-TEXT}}
\end{changes}
{{POINTER-TO-WHERE-IT-LANDED-AND-WHAT-IT-SETTLES}}
\end{reviewerresponse}


% Add further reviewers here, e.g.:
% \input{Reviewers/R2.tex}
% \input{Reviewers/R3.tex}

\end{document}
